% ============================================================================
%  SECCIÓN 6.5: MANEJO DE ESTADOS DE CARGA, ERRORES Y PRUEBAS
% ============================================================================

\section{Manejo de estados de carga, errores y pruebas}

Una aplicaci\'on conectada a un servidor depende de la red. Por ello es
necesario contemplar los momentos de espera mientras se obtienen los
datos y los errores que pueden ocurrir al comunicarse con Supabase,
de manera que la aplicaci\'on nunca se cierre de forma inesperada.

\subsection{Estados de carga}

Mientras la aplicaci\'on consulta Supabase existe un tiempo de espera.
Durante ese lapso la interfaz muestra un indicador de carga y, cuando los
datos llegan, se muestran en la lista:

\begin{figure}[H]
\centering
  \figuraAPA{fig:estados-carga}{Manejo del estado de carga en la interfaz}{%
  \begin{tikzpicture}[
    marco/.style={rectangle, draw=black!60, rounded corners,
      minimum width=4.4cm, minimum height=5.6cm,
      align=center, font=\footnotesize},
    fila/.style={rectangle, draw=black!30, minimum width=3.6cm,
      minimum height=0.55cm, align=center, font=\footnotesize},
    salto/.style={-Stealth, thick}
  ]
    \node[marco] (carga) {ESCENARIOS};
    \node[below=0.8cm of carga.north] {Cargando...};
    \node[below=0.5cm of carga.center] {\textcolor{gray}{\Large \ldots}};
    \node[below=0.12cm of carga.center] {\textcolor{gray}{\Large \ldots}};
    \node[above=0.5cm of carga.south] {\textcolor{gray}{\Large \ldots}};

    \node[marco, right=1.4cm of carga] (datos) {ESCENARIOS};
    \node[fila, below=0.9cm of datos.north] (f1) {Estadio Hernando Siles};
    \node[fila, below=0.12cm of f1] (f2) {Coliseo Eduardo Leon};
    \node[fila, below=0.12cm of f2] (f3) {Complejo deportivo Capriles};

    \draw[salto] ($(carga.east)+(0.15,0)$) -- node[above,
      font=\footnotesize] {carga} ($(datos.west)+(-0.15,0)$);
  \end{tikzpicture}
  }{Elaboraci\'on propia.}
\end{figure}

En la implementaci\'on, React Query expone autom\'aticamente los estados
de la consulta. La interfaz decide qu\'e mostrar seg\'un el estado:

\begin{lstlisting}
const { data, isLoading, isError, refetch } = useScenarios();

if (isLoading) {
  return <LoadingSpinner />;
}

if (isError) {
  return <EmptyState message="No se pudo obtener la informacion" />;
}

// datos disponibles: renderizar la lista
return <FlatList data={data} ... />;
\end{lstlisting}

\subsection{Manejo de errores de conexi\'on}

La aplicaci\'on contempla al menos los siguientes escenarios de error:

\begin{itemize}
  \item \textbf{Sin conexi\'on:} no se pudo obtener la informaci\'on del
    servidor. Se muestra un mensaje claro y la opci\'on de reintentar.
  \item \textbf{Error al guardar:} fall\'o la operaci\'on de escritura.
    Se informa al usuario para que intente nuevamente.
  \item \textbf{Datos inv\'alidos:} las validaciones de formulario (Zod)
    evitan enviar datos incorrectos al servidor.
\end{itemize}

Tanto las consultas como las mutaciones capturan el error devuelto por
Supabase y lo transforman en un mensaje legible. Por ejemplo, en el hook
de consulta:

\begin{lstlisting}
async function fetchScenarios() {
  const { data, error } = await supabase
    .from('scenarios')
    .select('*');

  if (error) {
    throw new Error(error.message); // lo captura React Query
  }
  return data;
}
\end{lstlisting}

De esta forma, si el servidor no est\'a disponible, la aplicaci\'on
muestra un estado de error controlado en lugar de cerrarse.

\subsection{Pruebas del CRUD}

Se realizaron pruebas de las cuatro operaciones y del manejo de errores
para verificar el comportamiento de la integraci\'on:

\begin{table}[H]
  \centering
  \caption{Pruebas realizadas sobre el CRUD con Supabase}
  \contenidoConFuente{%
    \begin{tablaAPA}
      \begin{tabular}{@{}p{2.2cm}p{4.4cm}p{6.6cm}@{}}
        \toprule
        \textbf{Prueba} & \textbf{Procedimiento} &
        \textbf{Resultado esperado} \\
        \midrule
        Prueba 1: CREATE &
        Crear un escenario desde el formulario &
        El registro aparece en la tabla de Supabase. \\
        \addlinespace
        Prueba 2: READ &
        Abrir la aplicaci\'on y consultar escenarios &
        La lista muestra datos de la base de datos. \\
        \addlinespace
        Prueba 3: UPDATE &
        Editar la capacidad de un escenario &
        El cambio se refleja al volver a consultar. \\
        \addlinespace
        Prueba 4: DELETE &
        Eliminar un evento o favorito &
        El registro desaparece de Supabase. \\
        \addlinespace
        Prueba 5: ERROR &
        Desconectar la red e intentar una consulta &
        Se muestra el mensaje de error controlado. \\
        \bottomrule
      \end{tabular}
    \end{tablaAPA}%
  }{Elaboraci\'on propia en base a las pruebas ejecutadas.}
\end{table}

El resultado de las pruebas confirma que las operaciones leen y escriben
correctamente sobre la base de datos centralizada y que la interfaz
responde de forma adecuada ante los estados de carga y error.